% !TEX TS-program = lualatex
\documentclass[11pt]{article}
\usepackage{amsmath,amssymb,amsthm}
\usepackage[margin=2.5cm]{geometry}
\usepackage{hyperref}
\usepackage{luatexja}

\theoremstyle{definition}
\newtheorem{theorem}{定理}
\newtheorem{definition}[theorem]{定義}

\title{フーリエ解析とその応用}
\author{田中 太郎}
\date{2026年3月}

\begin{document}
\maketitle

\section{フーリエ級数}

周期 $2\pi$ の関数 $f(x)$ は、次のフーリエ級数で展開できる：

\begin{equation}
  f(x) = \frac{a_0}{2} + \sum_{n=1}^{\infty}
  \left( a_n \cos nx + b_n \sin nx \right)
\end{equation}

ここでフーリエ係数は次式で与えられる：
\begin{align}
  a_n & = \frac{1}{\pi} \int_{-\pi}^{\pi} f(x) \cos nx \, dx
  \quad (n = 0, 1, 2, \ldots)                                \\
  b_n & = \frac{1}{\pi} \int_{-\pi}^{\pi} f(x) \sin nx \, dx
  \quad (n = 1, 2, 3, \ldots)
\end{align}

\section{フーリエ変換}

\begin{definition}[フーリエ変換]
  $f \in L^1(\mathbb{R})$ に対し、フーリエ変換 $\hat{f}$ を
  \begin{equation}
    \hat{f}(\xi) = \int_{-\infty}^{\infty}
    f(x)\, e^{-2\pi i x \xi} \, dx
  \end{equation}
  で定義する。
\end{definition}

\begin{theorem}[パーセバルの等式]
  $f \in L^2(\mathbb{R})$ ならば
  \begin{equation}
    \int_{-\infty}^{\infty} |f(x)|^2 \, dx
    = \int_{-\infty}^{\infty} |\hat{f}(\xi)|^2 \, d\xi
  \end{equation}
  が成り立つ。すなわちフーリエ変換はエネルギーを保存する。
\end{theorem}

\section{応用：熱方程式}

一次元熱方程式
\begin{equation}
  \frac{\partial u}{\partial t}
  = \alpha \frac{\partial^2 u}{\partial x^2},
  \qquad x \in \mathbb{R},\; t > 0
\end{equation}
の初期条件 $u(x,0) = f(x)$ に対する解は、
フーリエ変換を用いて次のように表される：
\begin{equation}
  u(x,t) = \frac{1}{\sqrt{4\pi\alpha t}}
  \int_{-\infty}^{\infty}
  f(y)\, \exp\!\left( -\frac{(x-y)^2}{4\alpha t} \right) dy
\end{equation}

\end{document}

